\begin{houseviewbox}[核心判断]
通鼎互联不是没有业绩的题材股：2026Q1扣非净利润为CNY1.03亿元，7月14日公告的H1扣非净利润预告为CNY2.30--3.30亿元，收入同比增长56--66\%。但公司也不是已经证明为AI高纯度成长股：2025年光纤光缆收入仅CNY1.83亿元、占集团收入5.34\%，电力电缆、通信电缆和安全业务仍是主要底盘。基于当前20.86元，我们给出基准目标价9.16元，综合目标价10.09元，三情景区间5.07--14.96元。
\end{houseviewbox}

\section{价格、目标价与行动}
\begin{exhibitbox}[表1-1\quad 决策仪表盘]
\begin{tabularx}{\exhibitboxwidth}{L{3.0cm}X L{3.0cm}X}
\textbf{实时价格} & CNY20.86（2026-07-14收盘） & \textbf{综合目标价} & CNY10.09 \\
\textbf{基本面目标} & CNY9.16（PE/PB/PS） & \textbf{合理区间} & CNY5.07--14.96 \\
\textbf{现价隐含} & 基准2026E扣非PE约49.3x & \textbf{行动} & 事件后观察，回撤验证，不追高 \\
\textbf{H1预告} & 归母1.70--2.40亿元；扣非2.30--3.30亿元 & \textbf{证据质量} & 官方为主，券商Street权重0\% \\
\end{tabularx}
\sourcenote{data/current\_valuation\_model\_20260714.json；data/verified\_market\_data.md；sources/official-20260714/2026-07-15-2026年半年度业绩预告.pdf}
\end{exhibitbox}

当前价格较综合目标价高106.7\%，并不等于股价必然立即回落；它表示市场在为H1兑现、H2光纤量价延续、安全业务利润和AI/数据中心期权同时付费。AStock的判断是，H1公告足以把公司从“亏损底部”抬升到“业绩修复”，但尚不足以把一个综合线缆平台重分类为高纯度成长股。

\section{投资逻辑排序}
\begin{enumerate}
\item \textbf{已确认：}2026Q1扣非净利润同比大幅扭亏，H1官方预告进一步确认公司层面盈利修复。
\item \textbf{已确认但仍需持续：}公司称光纤产品需求量价齐升，2025年光纤销量和库存也出现显著变化。
\item \textbf{可选项而非分母：}运营商集采、算力基础设施和特种光纤需求提供行业背景，但没有直接转化为公司H1分部收入的公开桥。
\item \textbf{估值约束：}当前价对应基准2026E扣非利润约49.3倍PE；若现金流、毛利率或H1分部数据不能跟上，估值锚应向下收敛。
\end{enumerate}

\section{三情景与行为}
\begin{exhibitbox}[表1-2\quad 三情景目标价]
\begin{tabularx}{\exhibitboxwidth}{L{1.3cm}R{1.8cm}R{1.8cm}X L{2.7cm}}
\textbf{情景} & \textbf{扣非利润} & \textbf{目标价} & \textbf{关键条件} & \textbf{行为} \\
熊 & CNY3.40亿 & CNY5.07 & H1处于区间下沿，H2利润和光纤价格回落 & 降低风险预算 \\
基准 & CNY5.20亿 & CNY9.16 & H1中枢兑现，H2现金流和毛利温和改善 & 事件后观察、回撤验证 \\
牛 & CNY7.00亿 & CNY14.96 & H1上沿、H2高增长且现金回款同步改善 & 仅在硬证据确认后跟随 \\
\end{tabularx}
\sourcenote{analysis/valuation\_model.md；data/current\_valuation\_model\_20260714.json}
\end{exhibitbox}
